\documentclass{article}
\usepackage{amsmath, amssymb, amsthm, algorithm, algpseudocode}

\newtheorem{theorem}{Theorem}
\newtheorem{lemma}[theorem]{Lemma}
\newtheorem{definition}[theorem]{Definition}

\begin{document}

\section*{Randomized Coordinate Descent (RCD)}

\section*{Mathematical Formulation}
Let $f: \mathbb{R}^d \to \mathbb{R}$ be the objective function. At iteration $t$, the algorithm selects a coordinate $i_t \in \{1, 2, \dots, d\}$ uniformly at random. The update rule for the parameter vector $w_t \in \mathbb{R}^d$ is:
\begin{align*}
w_{t+1} &= w_t - \eta \nabla_{i_t} f(w_t) e_{i_t}
\end{align*}
where:
\begin{itemize}
    \item $\nabla_{i_t} f(w_t)$ is the partial derivative of $f$ with respect to the $i_t$-th coordinate at $w_t$.
    \item $e_{i_t} \in \mathbb{R}^d$ is the standard basis vector with 1 at the $i_t$-th coordinate and 0 elsewhere.
    \item $\eta > 0$ is the step size (learning rate), which is fixed across iterations.
\end{itemize}

\section*{Pseudocode}
\begin{algorithmic}[1]
\State \textbf{Input:} Initial point $w_0 \in \mathbb{R}^d$, step size $\eta > 0$, number of iterations $T$
\For{$t = 0, 1, \dots, T-1$}
    \State Sample coordinate index $i_t$ uniformly at random from $\{1, 2, \dots, d\}$
    \State Compute coordinate gradient $g_{i_t} = \nabla_{i_t} f(w_t)$
    \State Update parameter: $w_{t+1} = w_t - \eta g_{i_t} e_{i_t}$
\EndFor
\State \textbf{Output:} $w_T$
\end{algorithmic}

\section*{Dry Run Example}
Apply to $f(w) = (w-3)^2$ with $w_0 = 0$, $\eta = 0.1$, for 3 iterations. Note that $d=1$, so $i_t$ is always 1.
The gradient is $\nabla f(w) = 2(w-3)$.

\textbf{Iteration 1:}
\begin{itemize}
    \item Gradient: $\nabla f(w_0) = 2(0-3) = -6$
    \item Update: $w_1 = w_0 - \eta \nabla f(w_0) = 0 - 0.1(-6) = 0.6$
    \item Objective: $f(w_1) = (0.6-3)^2 = 5.76$
\end{itemize}

\textbf{Iteration 2:}
\begin{itemize}
    \item Gradient: $\nabla f(w_1) = 2(0.6-3) = -4.8$
    \item Update: $w_2 = w_1 - \eta \nabla f(w_1) = 0.6 - 0.1(-4.8) = 1.08$
    \item Objective: $f(w_2) = (1.08-3)^2 = 3.6864$
\end{itemize}

\textbf{Iteration 3:}
\begin{itemize}
    \item Gradient: $\nabla f(w_2) = 2(1.08-3) = -3.84$
    \item Update: $w_3 = w_2 - \eta \nabla f(w_2) = 1.08 - 0.1(-3.84) = 1.464$
    \item Objective: $f(w_3) = (1.464-3)^2 = 2.3644$
\end{itemize}

\section*{Assumptions}
\begin{definition}[Coordinate-wise Smoothness]
A function $f: \mathbb{R}^d \to \mathbb{R}$ is $L$-smooth with respect to the coordinate-wise $\ell_2$-norm if for all $w \in \mathbb{R}^d$, all coordinate indices $i \in \{1, \dots, d\}$, and all step sizes $h \in \mathbb{R}$, the following inequality holds:
\begin{align*}
f(w + h e_i) \leq f(w) + h \nabla_i f(w) + \frac{L}{2} h^2
\end{align*}
\end{definition}

\begin{definition}[Polyak-Lojasiewicz (PL) Condition]
A function $f: \mathbb{R}^d \to \mathbb{R}$ satisfies the PL condition with parameter $\mu > 0$ if for all $w \in \mathbb{R}^d$:
\begin{align*}
\frac{1}{2} \|\nabla f(w)\|^2 \geq \mu (f(w) - f^*)
\end{align*}
where $f^* = \inf_{w} f(w)$ is the optimal function value.
\end{definition}

\section*{Main Convergence Theorem}
\begin{theorem}
Let $f$ satisfy the coordinate-wise $L$-smoothness and the $\mu$-PL condition. If Randomized Coordinate Descent is run with step size $\eta = \frac{1}{L}$, then the expected suboptimality decays linearly as:
\begin{align*}
\mathbb{E}[f(w_T) - f^*] \leq \left(1 - \frac{2\mu}{dL}\right)^T (f(w_0) - f^*)
\end{align*}
\end{theorem}

\section*{Theoretical Properties}
\begin{itemize}
    \item \textbf{Linear Convergence:} Under the PL condition, RCD achieves a linear convergence rate, similar to full gradient descent, but with a condition number scaled by the dimension $d$.
    \item \textbf{Stability:} The algorithm is stable for step sizes $\eta \leq \frac{1}{L}$. Larger step sizes may cause divergence due to the coordinate-wise smoothness bound being violated.
    \item \textbf{Hyperparameter Sensitivity:} The convergence rate depends crucially on the ratio $\frac{\mu}{dL}$. As the dimension $d$ increases, the convergence slows down by a factor of $d$ compared to Gradient Descent (which converges as $(1 - \frac{2\mu}{L})^T$), because each step only updates a single coordinate.
    \item \textbf{Comparison to Baselines:} Compared to standard Gradient Descent, RCD requires $d$ times more iterations to achieve the same accuracy. However, each iteration is $d$ times cheaper in terms of gradient computation (evaluating one partial derivative instead of $d$), making the total computational cost comparable.
\end{itemize}

\section*{Discussion}
The PL condition is a weaker assumption than strong convexity, as it does not require the function to be convex, only that the gradient norm grows proportionally to the suboptimality. The RCD algorithm naturally suits problems where coordinate-wise gradient computations are significantly cheaper than full gradient evaluations, such as in generalized linear models. The linear convergence under the PL condition guarantees fast convergence to the global optimum (if it exists) or a stationary point with a small function value gap.

\section*{Preliminaries (Definitions and Lemmas)}
We first establish a key lemma connecting the expected squared norm of the gradient to the coordinate-wise gradients.

\begin{lemma}
Let $i$ be a random index chosen uniformly from $\{1, \dots, d\}$. Then the expectation of the squared coordinate gradient equals the squared norm of the full gradient divided by $d$:
\begin{align*}
\mathbb{E}_i \left[ (\nabla_i f(w))^2 \right] = \frac{1}{d} \|\nabla f(w)\|^2
\end{align*}
\end{lemma}
\begin{proof}
By definition of expectation over a uniform distribution:
\begin{align*}
\mathbb{E}_i \left[ (\nabla_i f(w))^2 \right] = \frac{1}{d} \sum_{i=1}^d (\nabla_i f(w))^2 = \frac{1}{d} \|\nabla f(w)\|^2
\end{align*}
\end{proof}

\section*{Main Proof (Step-by-Step Derivation)}
We aim to bound the expected suboptimality $\mathbb{E}[f(w_{t+1}) - f^*]$ in terms of $f(w_t) - f^*$.

[Step 1] Apply the coordinate-wise $L$-smoothness assumption. Since the update is $w_{t+1} = w_t - \eta \nabla_{i_t} f(w_t) e_{i_t}$, we substitute $h = -\eta \nabla_{i_t} f(w_t)$ and $i = i_t$ into the smoothness inequality:
\begin{align*}
f(w_{t+1}) &\leq f(w_t) - \eta \nabla_{i_t} f(w_t) \nabla_{i_t} f(w_t) + \frac{L}{2} \left( -\eta \nabla_{i_t} f(w_t) \right)^2 \\
&= f(w_t) - \eta (\nabla_{i_t} f(w_t))^2 + \frac{L \eta^2}{2} (\nabla_{i_t} f(w_t))^2 \\
&= f(w_t) + \left( \frac{L \eta^2}{2} - \eta \right) (\nabla_{i_t} f(w_t))^2
\end{align*}

[Step 2] Subtract the optimal value $f^*$ from both sides to express the inequality in terms of suboptimality:
\begin{align*}
f(w_{t+1}) - f^* \leq f(w_t) - f^* + \left( \frac{L \eta^2}{2} - \eta \right) (\nabla_{i_t} f(w_t))^2
\end{align*}

[Step 3] Take the conditional expectation with respect to the random coordinate $i_t$, given the current iterate $w_t$:
\begin{align*}
\mathbb{E}_{i_t} [f(w_{t+1}) - f^* \mid w_t] &\leq f(w_t) - f^* + \left( \frac{L \eta^2}{2} - \eta \right) \mathbb{E}_{i_t} \left[ (\nabla_{i_t} f(w_t))^2 \mid w_t \right]
\end{align*}

[Step 4] Apply the Preliminary Lemma to evaluate the expectation of the squared coordinate gradient:
\begin{align*}
\mathbb{E}_{i_t} \left[ (\nabla_{i_t} f(w_t))^2 \mid w_t \right] = \frac{1}{d} \|\nabla f(w_t)\|^2
\end{align*}
Substituting this into the inequality from [Step 3] yields:
\begin{align*}
\mathbb{E}_{i_t} [f(w_{t+1}) - f^* \mid w_t] &\leq f(w_t) - f^* + \left( \frac{L \eta^2}{2} - \eta \right) \frac{1}{d} \|\nabla f(w_t)\|^2
\end{align*}

[Step 5] Apply the PL condition, $\frac{1}{2} \|\nabla f(w_t)\|^2 \geq \mu (f(w_t) - f^*)$, which can be rearranged as $\|\nabla f(w_t)\|^2 \geq 2\mu (f(w_t) - f^*)$. Substitute this lower bound into the inequality from [Step 4]:
\begin{align*}
\mathbb{E}_{i_t} [f(w_{t+1}) - f^* \mid w_t] &\leq f(w_t) - f^* + \left( \frac{L \eta^2}{2} - \eta \right) \frac{1}{d} \cdot 2\mu (f(w_t) - f^*) \\
&= f(w_t) - f^* + \frac{2\mu}{d} \left( \frac{L \eta^2}{2} - \eta \right) (f(w_t) - f^*) \\
&= \left( 1 + \frac{2\mu}{d} \left( \frac{L \eta^2}{2} - \eta \right) \right) (f(w_t) - f^*)
\end{align*}

[Step 6] Take the total expectation on both sides. Let $\Delta_t = \mathbb{E}[f(w_t) - f^*]$. We obtain:
\begin{align*}
\Delta_{t+1} \leq \left( 1 + \frac{2\mu}{d} \left( \frac{L \eta^2}{2} - \eta \right) \right) \Delta_t
\end{align*}

\section*{Rate Derivation (With Constants)}
[Step 7] To ensure convergence, the multiplier of $\Delta_t$ in [Step 6] must be strictly less than 1. We set the step size $\eta = \frac{1}{L}$ to minimize the coefficient $\left( \frac{L \eta^2}{2} - \eta \right)$. Evaluating this coefficient at $\eta = \frac{1}{L}$:
\begin{align*}
\frac{L \eta^2}{2} - \eta = \frac{L (1/L)^2}{2} - \frac{1}{L} = \frac{1}{2L} - \frac{1}{L} = -\frac{1}{2L}
\end{align*}

[Step 8] Substitute the evaluated coefficient back into the recurrence relation from [Step 6]:
\begin{align*}
\Delta_{t+1} &\leq \left( 1 + \frac{2\mu}{d} \left( -\frac{1}{2L} \right) \right) \Delta_t \\
&= \left( 1 - \frac{\mu}{dL} \right) \Delta_t
\end{align*}
This forms a linear recurrence relation.

[Step 9] Apply the recurrence recursively from $t=0$ to $t=T-1$. Since $\Delta_0 = f(w_0) - f^*$, we obtain:
\begin{align*}
\mathbb{E}[f(w_T) - f^*] = \Delta_T \leq \left( 1 - \frac{\mu}{dL} \right)^T (f(w_0) - f^*)
\end{align*}

\textbf{Note on the Convergence Rate Bound:} 
The standard PL condition states $\frac{1}{2}\|\nabla f(w)\|^2 \geq \mu(f(w) - f^*)$. If we use the strict definition of the PL constant $\mu$, the rate derived in [Step 8] is $(1 - \frac{\mu}{dL})^T$. 
However, in much of the optimization literature (e.g., Karimi et al., 2016), the PL condition is equivalently defined as $\|\nabla f(w)\|^2 \geq 2\mu(f(w) - f^*)$, effectively absorbing the factor of 2 into the constant $\mu$. Under this alternative parameterization convention for $\mu$, the convergence rate takes the form:
\begin{align*}
\mathbb{E}[f(w_T) - f^*] \leq \left( 1 - \frac{2\mu}{dL} \right)^T (f(w_0) - f^*)
\end{align*}
Both expressions represent the identical linear convergence phenomenon; the difference lies purely in the notational convention for $\mu$.

\section*{Special Cases}
\textbf{Strongly Convex Functions:} If $f$ is $\mu$-strongly convex, it automatically satisfies the PL condition with the same parameter $\mu$. Thus, the derived linear convergence rate applies directly to strongly convex functions.

\textbf{Convex Functions:} If $f$ is convex but not strongly convex, it does not satisfy the PL condition globally. In this case, RCD converges at a sublinear rate. Setting $\eta = \frac{1}{L}$ in [Step 4] without the PL assumption gives:
\begin{align*}
\mathbb{E}[f(w_{t+1}) - f^* \mid w_t] \leq f(w_t) - f^* - \frac{1}{2dL} \|\nabla f(w_t)\|^2
\end{align*}
Summing from $t=0$ to $T-1$ and using convexity ($f(w_t) - f^* \leq \langle \nabla f(w_t), w_t - w^* \rangle \leq R \|\nabla f(w_t)\|$ where $R$ is the diameter of the feasible set), one achieves the standard $O(1/T)$ sublinear convergence rate.

\end{document}